\documentclass{beamer}

\usepackage{listings,tikz}
\usetikzlibrary{calc,arrows,decorations.markings,positioning}

% Setup epigraph.
\usepackage{epigraph,etoolbox}
\setlength\epigraphwidth{10cm}
\setlength\epigraphrule{0.5pt}
\makeatletter
\patchcmd{\epigraph}{\@epitext{#1}}{\itshape\@epitext{#1}}{}{}
\makeatother

\tikzset{
  nrm/.style={
    draw=none,
    rectangle,
    rounded corners=2mm,
    inner sep=3mm,
    fill=Base02,
    ultra thick,
    minimum width=5cm,
    text=Base04
  },
  lit/.style={
    draw=Base0D,
    rectangle,
    rounded corners=2mm,
    inner sep=3mm,
    fill=Base0D,
    ultra thick,
    minimum width=5cm,
    text=Base07
  },
  bad/.style={
    draw=Base08,
    rectangle,
    rounded corners=2mm,
    inner sep=3mm,
    fill=Base08,
    ultra thick,
    minimum width=5cm,
    text=Base07
  },
  exp/.style={
    draw=Base0D,
    rectangle,
    rounded corners=2mm,
    inner sep=3mm,
    fill=Base02,
    ultra thick,
    text width=4cm,
    text=Base07,
    text depth=4.7cm
  },
  arr/.style={
    draw=Base0D,
    fill=Base0D,
    line width=2pt,
    dashed
    %% decoration={markings,mark=at position 1 with {\arrow[scale=2,thin,Base0D]{triangle 45}}},
    %% shorten >= 8pt,
    %% preaction={decorate},
    %% postaction={draw,line width=5pt}
  },
  bb/.style={
    rectangle,
    draw=none,
    fill=Base00,
    inner sep=3mm,
    ultra thick,
    rounded corners=1mm,
    text width=3cm,
    text=Base04,
    align=center
  },
  wb/.style={
    rectangle,
    draw=none,
    fill=Base06,
    inner sep=3mm,
    ultra thick,
    rounded corners=1mm,
    text width=3cm,
    text=Base02,
    align=center
  },
  gb/.style={
    rectangle,
    draw=none,
    fill=Base03,
    inner sep=3mm,
    ultra thick,
    rounded corners=1mm,
    text width=3cm,
    text=Base06,
    align=center
  }
}

%% \lstset{
\lstdefinestyle{C}{
  language=C++,
  %% numbers=left,
  showstringspaces=false,
  formfeed=\newpage,
  tabsize=4,
  basicstyle=\footnotesize\ttfamily,
  commentstyle=\color{Base08}\itshape,
  keywordstyle=\color{Base0D},
  stringstyle=\color{Base0B},
  morekeywords={std, thrust, include, ifndef, define, endif}
  %% morekeywords={models, lambda, forms, dict, list, str, import, dir, help,
  %%  zip, with, open}
}
%% \lstset{
\lstdefinestyle{Py}{
  language=Python,
  %% numbers=left,
  showstringspaces=false,
  formfeed=\newpage,
  tabsize=4,
  basicstyle=\footnotesize\ttfamily,
  commentstyle=\color{Base08}\itshape,
  keywordstyle=\color{Base0D},
  stringstyle=\color{Base0B},
  morekeywords={assert}
  %% morekeywords={models, lambda, forms, dict, list, str, import, dir, help,
  %%  zip, with, open}
}
%% \lstdefinestyle{C}{language=C}
%% \lstdefinestyle{Py}{language=Python}

\mode<presentation>{\usetheme{Median16}}
\title[Software Testing for Researchers]{Software Testing for Researchers}
\author{Luke Hodkinson}
\institute{\hfill Center for Astrophysics and Supercomputing \\
  \hfill Swinburne University of Technology \\
  \hfill Melbourne, Hawthorn 32000, \underline{Australia} }

\date{\today}
\titlegraphic{
  \begin{tikzpicture}[overlay]
    \node[anchor=south west,inner sep=0pt] at ($(current page.south west)+(7cm,2cm)$) {
      \includegraphics[width=3cm]{swinburne.png}
    };
  \end{tikzpicture}
  \vspace{-\baselineskip}
}

\begin{document}

\frame{\titlepage}

%% \begin{frame}
  \frametitle{Motivation}
  Are any of these phrases familiar?
  \begin{itemize}
    \pause
    \item It was working just last week!
    \pause
    \item I only made a small change, I don't know how it could have broken \emph{that}...
    \pause
    \item I'm sure I've fixed this before.
  \end{itemize}
\end{frame}

\begin{frame}
  \frametitle{Motivation}
  Tests will save you time.
  \begin{itemize}
    \item Larger codes will have many systems.
    \item Unexpected consequences.
    \item Automation.
    \item Less debugging later.
  \end{itemize}
\end{frame}

\begin{frame}
  \frametitle{Motivation}
  Show stats on typical code sizes and sizes of other test suites to highlight
  time savings. Also show complexity with stats on if conditionals.
\end{frame}

\begin{frame}
  \frametitle{Motivation}
  Tests don't just idenitfy problems, they prevent them.

  What here?
\end{frame}

\begin{frame}[fragile]
  \frametitle{Motivation}
  Tests make your code more attractive.
  \epigraph{``Code without tests is broken by design.''}{--- \textup{Jacob Kaplan-Moss}, Django documentation}
\end{frame}

\begin{frame}
  \frametitle{Motivation}
  Tests help teams work together.

  If you modify someone else's code, how do you know it all still works?
\end{frame}

\begin{frame}
  \frametitle{Motivation}
  Show massive list of projects with test suites.
\end{frame}


\begin{frame}
  \frametitle{Testing Methods}
  \begin{tikzpicture}[overlay]

    \node[bkg,anchor=north west,text width=3.2cm,text depth=6cm,label=above:{\color{Base04}Static}] at (-0.5,2.3) {};
    \node[bkg,anchor=north west,text width=6.8cm,text depth=6cm,label=above:{\color{Base04}Dynamic}] at (3.8,2.3) {};

    \visible<2->{

    \node[nrm2,short,anchor=north west] at (-0.2,1.9) {Reviews};
    \node[nrm2,short,anchor=north west] at (-0.2,0.7) {Walkthroughs};
    \node[nrm2,short,anchor=north west] at (-0.2,-0.5) {Inspections};
    \node[nrm2,short,anchor=north west] at (-0.2,-1.7) {IDE (syntax)};

    }

    \visible<2-4>{
      \node[nrm2,short,anchor=north west] at (-0.2,-2.9) {Analysis};
    }

    \visible<3->{

    \node[nrm2,long,anchor=north west] (F) at (4.1,1.9) {Functional};
    \node[nrm2,long,anchor=north west] (N) at (4.1,-1.6) {Non-functional};

    }

    \coordinate (o) at (6,0.7);
    \coordinate (p) at (6,-2.8);
    \coordinate (g) at (7.52,-3);

    \visible<4>{

      \node[nrm2,tiny,anchor=north west] (a) at (4.1,0.5)  {Unit};
      \node[nrm2,tiny,anchor=north west] (b) at (5.1,-0.4) {Integration};
      \node[nrm2,tiny,anchor=north west] (c) at (6.72,0.5)  {System};
      \node[nrm2,tiny,anchor=north west] (d) at (7.8,-0.4) {Regression};

    }

    \visible<5>{

    \node[lit,tiny,anchor=north west] (a) at (4.1,0.5)  {Unit};
    \node[lit,tiny,anchor=north west] (b) at (5.1,-0.4) {Integration};
    \node[lit,tiny,anchor=north west] (c) at (6.72,0.5)  {System};
    \node[lit,tiny,anchor=north west] (d) at (7.8,-0.4) {Regression};
    \node[lit,short,anchor=north west] at (-0.2,-2.9) {Analysis};
    \node[lit,tiny,anchor=north west] (e) at (9.4,0.5)  {Sanity};

    }

    \visible<4>{
      \node[nrm2,tiny,anchor=north west] (e) at (9.4,0.5)  {Sanity};
    }

    \visible<4->{

    \node[nrm2,tiny,anchor=north west] (f) at (4.1,-3)  {Performance};
    \node[nrm2,tiny,anchor=north west]     at (6.95,-3) {Usability};
    \node[nrm2,tiny,anchor=north west] (h) at (9.18,-3)  {Security};

    \path[draw,color=Base06,-latex] (F) |- (o) -| (a);
    \path[draw,color=Base06,-latex] (F) |- (o) -| (b);
    \path[draw,color=Base06,-latex] (F) |- (o) -| (c);
    \path[draw,color=Base06,-latex] (F) |- (o) -| (d);
    \path[draw,color=Base06,-latex] (F) |- (o) -| (e);

    \path[draw,color=Base06,-latex] (N) |- (p) -| (f);
    \path[draw,color=Base06,-latex] (N) |- (p) -| (g);
    \path[draw,color=Base06,-latex] (N) |- (p) -| (h);

    }

  \end{tikzpicture}
\end{frame}

\begin{frame}
  \frametitle{The ``Box'' Approach}
  Describes the point-of-view of the tester.
  Divided into three boxes:
  \begin{itemize}
  \item Black-box testing.
  \item White-box testing.
  \item Grey-box testing.
  \end{itemize}
  Why is this useful? Helps decide on ecapsulation and
  implementation hiding, and with TDD.
\end{frame}

\begin{frame}
  \frametitle{Black-box Testing}
  We don't know anything about internals, only
  the interfaces.

  TODO: Include picture.
\end{frame}

\begin{frame}[fragile]
  \frametitle{Black-box Testing}
  \begin{block}{Function definition}
    \begin{lstlisting}[style=C]
  // Perform Newton-Raphson solve.
  // Solution is bounded by [x1,x2].
  // Initial value is x.
  // Solve tolerance is tol, maximum iterations are
  // max_its.
  template< class FuncT, class T >
  T newton( FuncT func, T x1, T x2, T x,
            T tol, int max_its );
    \end{lstlisting}
  \end{block}
\end{frame}

\begin{frame}[fragile]
  \frametitle{Black-box Testing}
  \begin{block}{Potential tests}
    \begin{itemize}
    \item Check known functions work ($x$, $x^2$, etc).
    \item Check violation of lower/upper bounds.
    \item Check variation of initial x.
    \item Check violation of tolerance.
    \item Check violation of maximum iterations.
    \end{itemize}
  \end{block}
  \pause
  \vspace{1cm}
  {\LARGE\color{Base0A} Can we do better?}
\end{frame}

\begin{frame}
  \frametitle{White-box Testing}
  Have access to application internals.

  TODO: Include picture.
\end{frame}

\begin{frame}[fragile]
  \frametitle{White-box Testing}
  \begin{block}{Function implementation}
    \begin{lstlisting}[style=C]
  template< class FuncT, class T >
  T newton( FuncT func, T x1, T x2, T x,
            T tol, int max_its ) {
    ...
    // Use a scaling to make sure we don't
    // converge simply because x is very small.
    T scale;
    if( x > 0.0 )
      scale = 1.0/x;
    else
      scale = 1.0;
    ...
  }
    \end{lstlisting}
  \end{block}
\end{frame}

\begin{frame}[fragile]
  \frametitle{White-box Testing}
  \begin{block}{More thorough tests}
    \begin{itemize}
    \item Check systems with small scales converge correctly.
    \item Check systems with large scales converge correctly.
    \item Check initial value of 0 works as expected.
    \end{itemize}
  \end{block}
  \pause
  \vspace{1cm}
  {\LARGE\color{Base0A} Perhaps too much effort?}
\end{frame}

\begin{frame}[fragile]
  \frametitle{Grey-box Testing}
  Somewhere in between.
\end{frame}


%% \begin{frame}
  \frametitle{Assertions}
  \begin{tikzpicture}[overlay]
    \node[lit,anchor=north west] (expr) at (2.5,2) {{\Large Is expression true?}};
    \node[bad,anchor=north west] (no) at (5.5,-2) {{\Large Abort}};
    \node[lit,anchor=north west] (yes) at (-0.5,-2) {{\Large Continue}};
    \node[below right=of expr] (nolabel) {No};
    \node[below left=of expr] (yeslabel) {Yes};
    \draw[line width=1pt,draw=Base06,->,>=latex] (expr.south) |- (3,0) -| (yes.north);
    \draw[line width=1pt,draw=Base06,->,>=latex] (expr.south) |- (5,0) -| (no.north);
  \end{tikzpicture}
\end{frame}

\begin{frame}
  \frametitle{Assertions}
  \begin{block}{Advantages}
    \begin{itemize}
    \item Very simple.
    \item Already provided by C/C++ and Python.
    \item Optimised out.
    \end{itemize}
  \end{block}
  \begin{block}{Disadvantages}
    \begin{itemize}
    \item Does not provide much information.
    \item Not a test ``suite'' as such.
    \end{itemize}
  \end{block}
\end{frame}

\begin{frame}[fragile]
  \frametitle{Assertions}
  \begin{example}
    \begin{lstlisting}[style=C]
  char* alloc_a_string( size_t size ) {
    char* p = (char*)malloc( size );

    // What if the allocation fails?
    // Probably a segfault.

    return p;
  }
    \end{lstlisting}
  \end{example}
\end{frame}

\begin{frame}[fragile]
  \frametitle{Assertions}
  \begin{example}
    \begin{lstlisting}[style=C]
  #include <assert.h>

  char* alloc_a_string( size_t size ) {
    char* p = (char*)malloc( size );
    assert( p ); // check allocation
    return p;
  }
    \end{lstlisting}
  \end{example}
\end{frame}

\begin{frame}[fragile]
  \frametitle{Assertions}
  \begin{example}
    \begin{lstlisting}[style=C]
  #ifndef NDEBUG
    #define assert( expr ) <something>
  #else
    #define assert( expr )
  #endif

  // In debug mode.
  gcc -c -o program program.c

  // In optimised mode.
  gcc -DNDEBUG -c -o program program.c
    \end{lstlisting}
  \end{example}
\end{frame}

\begin{frame}[fragile]
  \frametitle{Assertions}
  \begin{example}
    \begin{lstlisting}[style=Py]
  def do_some_asserts(*args, **kw):

    # Must have at least one positional.
    assert len(args) > 0

    # Must have a key of 'hello'.
    assert kw.has('hello')

    # Can also have messages.
    assert kw['hello'] == 'world', 'Wrong!'
    
    \end{lstlisting}
  \end{example}
\end{frame}


%% \begin{frame}
  \frametitle{Testing Levels}
  \begin{tikzpicture}[overlay,shift={(1.5,-2.75)}]
    \visible<1,3->{
      \node[nrm,anchor=north west] at (-1.5,5) {{\Large Unit testing}};
    }
    \visible<2>{
      \node[lit,anchor=north west] at (-1.5,5) {{\Large Unit testing}};
      \node[exp,anchor=north west] at (4.5,5)  {{\small Tests
          the lowest level of the code, focusing on each individual function.}};
      \draw[arr] (4.4,4.5) to (3.6,4.5);
      \node[bb,anchor=north west] at (5,3) {{\small Black-box}};
      \node[wb,anchor=north west] at (5,2) {{\small White-box}};
      \node[gb,anchor=north west] at (5,1) {{\small Grey-box}};
    }

    \visible<-2,4->{
      \node[nrm,anchor=north west] at (-1.5,3.5)   {{\Large Integration testing}};
    }
    \visible<3>{
      \node[lit,anchor=north west] at (-1.5,3.5)   {{\Large Integration testing}};
      \node[exp,anchor=north west] at (4.5,5) {{\small Tests
          the interoperability of separate code pieces.}};
      \draw[arr] (4.4,3) to (3.6,3);
      \node[bb,anchor=north west] at (5,3) {{\small Black-box}};
      \node[gb,anchor=north west] at (5,2) {{\small Grey-box}};
    }

    \visible<-3,5->{
      \node[nrm,anchor=north west] at (-1.5,2) {{\Large System testing}};
    }
    \visible<4>{
      \node[lit,anchor=north west] at (-1.5,2) {{\Large System testing}};
      \node[exp,anchor=north west] at (4.5,5) {{\small Tests the entire system
          working together.}};
      \draw[arr] (4.4,1.5) to (3.6,1.5);
      \node[bb,anchor=north west] at (5,3) {{\small Black-box}};
      \node[gb,anchor=north west] at (5,2) {{\small Grey-box}};
    }

    \visible<-4>{
      \node[nrm,anchor=north west] at (-1.5,0.5)   {{\Large Regression testing}};
    }
    \visible<5>{
      \node[lit,anchor=north west] at (-1.5,0.5)   {{\Large Regression testing}};
      \node[exp,anchor=north west] at (4.5,5) {{\small Ensures previous errors
          don't find their way back in.}};
      \draw[arr] (4.4,0) to (3.6,0);
      \node[bb,anchor=north west] at (5,3) {{\small Black-box}};
    }
  \end{tikzpicture}
\end{frame}

\begin{frame}
  \frametitle{Unit Testing}
  Lowest level of testing.
  Usually directly tests inputs and outputs of functions/methods.
\end{frame}

\begin{frame}
  \frametitle{Unit Testing}
  Things we would like from unit testing:
  \begin{itemize}
  \vspace{0.5cm}
  \item Easy to implement.
  \vspace{0.5cm}
  \item Descriptive failures.
  \vspace{0.5cm}
  \item Minimum boilerplate.
  \vspace{0.5cm}
  \item Flexibility.
  \end{itemize}
\end{frame}

\begin{frame}[fragile]
  \frametitle{Unit Testing}
  \begin{block}{Code to test}
    \begin{lstlisting}[style=C]
  // Complex number implementation.
  struct complex {
    double real;
    double imag;

    complex add( complex const& other ) const;
    complex sub( complex const& other ) const;
    complex div( complex const& other ) const;
    complex mul( complex const& other ) const;
  };
    \end{lstlisting}
  \end{block}
\end{frame}

\begin{frame}[fragile]
  \frametitle{Unit Testing}
  \begin{block}{Test program 1}
    \begin{lstlisting}[style=C]
  int main() {
    complex a{ 1, 0 }, b{ 0, 1 }, c{ 1, 1 };

    assert( a.add( a ) == 2 );
    assert( a.sub( a ) == 0 );
    assert( a.add( b ) == c );
    assert( a.mul( a ) == a );
    assert( c.div( a ) == c );
  }
    \end{lstlisting}
  \end{block}
\end{frame}

\begin{frame}[fragile]
  \frametitle{Unit Testing}
  Things we would like from unit testing:
  \begin{itemize}
  \vspace{0.5cm}
  \item Easy to implement.
  \vspace{0.5cm}
  \item Descriptive failures.
  \vspace{0.5cm}
  \item Minimum boilerplate.
  \vspace{0.5cm}
  \item Flexibility.
  \end{itemize}
  \begin{tikzpicture}[overlay]
    \node[Base0B,scale=2] at (5.25,4) {\checkmark};
    \node[Base08,scale=2] at (5.25,2.8) {$\times$};
    \node[Base0B,scale=2] at (5.25,1.8) {\checkmark};
    \node[Base08,scale=2] at (5.25,0.6) {$\times$};
  \end{tikzpicture}
\end{frame}

\begin{frame}[fragile]
  \frametitle{Unit Testing}
  \begin{block}{Test program 2}
    \begin{lstlisting}[style=C]
  int main() {
    complex a{ 1, 0 }, b{ 0, 1 }, c{ 1, 1 };
    complex c{ 1, 1 };

    if( a.add( a ) != 2 )
      (printf( "addition failed" ), abort());
    if( a.add( b ) != c )
      (printf( "subtraction failed" ), abort());
    if( a.mul( a ) != a )
      (printf( "multiplication failed" ), abort());
    if( c.div( a ) != c )
      (printf( "division failed" ), abort());
  }
    \end{lstlisting}
  \end{block}
\end{frame}

\begin{frame}[fragile]
  \frametitle{Unit Testing}
  Things we would like from unit testing:
  \begin{itemize}
  \vspace{0.5cm}
  \item Easy to implement.
  \vspace{0.5cm}
  \item Descriptive failures.
  \vspace{0.5cm}
  \item Minimum boilerplate.
  \vspace{0.5cm}
  \item Flexibility.
  \end{itemize}
  \begin{tikzpicture}[overlay]
    \node[Base0B,scale=2] at (5.25,4) {\checkmark};
    \node[Base0B,scale=2] at (5.25,2.8) {\checkmark};
    \node[Base08,scale=2] at (5.25,1.8) {$\times$};
    \node[Base08,scale=2] at (5.25,0.6) {$\times$};
  \end{tikzpicture}
\end{frame}

\begin{frame}[fragile]
  \frametitle{Unit Testing}
  We \emph{could} continue adding more, but, as always, someone
  has probably already done this.
\end{frame}

\begin{frame}[fragile]
  \frametitle{C/C++ Unit Testing Frameworks}
  \begin{tabular}{p{3cm}|p{7cm}}
    {\color{Base09}CppUnit} & Excellent framework, very well used/tested, multiple loggers, many test assertions. \\
    \hline
    {\color{Base09}MinUnit} & Possibly the simplest framework available. \\
    \hline
    {\color{Base09}CxxTest} & Essentially no boilerplate, uses Python scripts to detect tests. \\
    \hline
    \only<1>{
      {\color{Base09}Catch} & Personal favorite, combines clever code and rich features.
    }
    \only<2>{
      {\color{Base09}Catch}\hspace{1cm}{\color{Base0B}$\checkmark$} & Personal favorite, combines clever code and rich features.
    }
  \end{tabular}
\end{frame}

\begin{frame}[fragile]
  \frametitle{Unit Testing}
  \begin{block}{Test program 3}
    \begin{lstlisting}[style=C]
  #include <catch.hpp>

  TEST_CASE( "Complex numbers" ) {
    complex a{ 1, 0 }, b{ 0, 1 }, c{ 1, 1 };
    REQUIRE( a.add( a ) == 2 );
    REQUIRE( a.sub( a ) == 0 );
    REQUIRE( a.add( b ) == c );
    REQUIRE( a.mul( a ) == a );
    REQUIRE( c.div( a ) == c );
  }
    \end{lstlisting}
  \end{block}
  \visible<2>{ \tikz[overlay] \draw[Base08,line width=2pt] (0.7,1.2) rectangle (2.3,3.3); }
\end{frame}

\begin{frame}[fragile]
  \frametitle{Successful Results}
  \begin{block}{Minimal output}
    All tests passed (5 assertions in 1 test case)
  \end{block}
  \begin{block}{Verbose output}
    \begin{verbatim}
  complex.cc:55: 
  PASSED:
    REQUIRE( c.div( a ) == c )
  with expansion:
    1 + 1i == 1 + 1i
    \end{verbatim}
  \end{block}
  \visible<2>{
  \begin{tikzpicture}[overlay]
    \draw[Base08,line width=2pt] (2.5,1.7) ellipse (2 and 0.5);
  \end{tikzpicture}
  }
\end{frame}

\begin{frame}
  \frametitle{Test Assertions}
  From {\color{Base0D} CxxTest}:
  \begin{columns}[onlytextwidth]
    \begin{column}{0.4\textwidth}
      \begin{itemize}
      \item {\ttfamily\small TS\_ASSERT}
      \item {\ttfamily\small TS\_ASSERT\_EQUALS}
      \item {\ttfamily\small TS\_ASSERT\_DIFFERS}
      \item {\ttfamily\small TS\_ASSERT\_DELTA}
      \end{itemize}
    \end{column}
    \begin{column}{0.6\textwidth}
      \begin{itemize}
      \item {\ttfamily\small TS\_ASSERT\_LESS\_THAN}
      \item {\ttfamily\small TS\_ASSERT\_LESS\_THAN\_EQUALS}
      \item {\ttfamily\small TS\_ASSERT\_GREATER\_EQUALS}
      \item {\ttfamily\small TS\_ASSERT\_GREATER\_THAN\_EQUALS}
      \end{itemize}
    \end{column}
  \end{columns}
  \vspace{0.5cm}
  From {\color{Base0D} Catch}:
  \begin{itemize}
  \item {\ttfamily\small REQUIRE}
  \end{itemize}
\end{frame}

\begin{frame}[fragile]
  \frametitle{Failed Results}
  \begin{verbatim}
-----------------------------------------------
Complex numbers
-----------------------------------------------
complex.cc:53
...............................................

complex.cc:60: FAILED:
  REQUIRE( c.div( a ) != c )
with expansion:
  1 + 1i != 1 + 1i

===============================================
1 test case - failed (5 assertions - 1 failed)
  \end{verbatim}
\end{frame}

\begin{frame}[fragile]
  \frametitle{Reporters}
  \footnotesize
  \begin{verbatim}
<Catch name="complex">
  <Group>
    <TestCase name="Complex numbers">
      <Expression success="false" filename="complex.cc" line="60">
        <Original>
          c.div( a ) != c
        </Original>
        <Expanded>
          1 + 1i != 1 + 1i
        </Expanded>
      </Expression>
      <OverallResult success="false"/>
    </TestCase>
    <OverallResults successes="4" failures="1"/>
  </Group>
  <OverallResults successes="4" failures="1"/>
</Catch>
  \end{verbatim}
\end{frame}

\begin{frame}[fragile]
  \frametitle{Unit Testing with ``Catch''}
  Things we would like from unit testing:
  \begin{itemize}
  \vspace{0.5cm}
  \item Easy to implement.
  \vspace{0.5cm}
  \item Descriptive failures.
  \vspace{0.5cm}
  \item Minimum boilerplate.
  \vspace{0.5cm}
  \item Flexibility.
  \end{itemize}
  \begin{tikzpicture}[overlay]
    \node[Base0B,scale=2] at (5.25,4) {\checkmark};
    \node[Base0B,scale=2] at (5.25,2.8) {\checkmark};
    \node[Base0B,scale=2] at (5.25,1.8) {\checkmark};
    \node[Base0B,scale=2] at (5.25,0.6) {\checkmark};
  \end{tikzpicture}
\end{frame}

\begin{frame}
  \frametitle{Fixtures}
\end{frame}

\begin{frame}
  \frametitle{Integration Testing}
\end{frame}

\begin{frame}
  \frametitle{System Testing}
\end{frame}

\begin{frame}
  \frametitle{Regression Testing}
\end{frame}

\begin{frame}
  \frametitle{Other Tests}
\end{frame}


\end{document}
